% Vietnamese translation for OI Public Library
% Source-File: IOI中国国家候选队论文/2006/陈启峰/附件/骑士解题报告.doc
\documentclass[11pt]{article}
\usepackage{oipl-vn}

\title{Lời giải bài ``Kỵ sĩ''}
\author{Chen Qifeng}
\date{}

\begin{document}
\maketitle

\origtitle{Knight solution report}
\sourcepath{IOI 2006 / Chen Qifeng / attachments / Knight solution report.doc}

\section{Mô tả bài toán}

Một quân kỵ sĩ di chuyển trên bàn cờ vô hạn. Mỗi kiểu di chuyển được mô tả
bằng một cặp số nguyên \((a,b)\): từ vị trí \((x,y)\), kỵ sĩ có thể nhảy đến
\((x+a,y+b)\) hoặc \((x-a,y-b)\). Mỗi kỵ sĩ có một danh sách hữu hạn các cặp
số nguyên, mô tả tất cả các kiểu nhảy mà nó được phép dùng. Bài toán bảo đảm
các vị trí mà kỵ sĩ có thể đến không cùng nằm trên một đường thẳng.

Hai kỵ sĩ được gọi là tương đương nếu, bắt đầu từ \((0,0)\), tập tất cả các
vị trí mà chúng có thể đến là hoàn toàn giống nhau, sau một số tùy ý bước
nhảy. Có thể chứng minh rằng với mọi kỵ sĩ đều tồn tại một kỵ sĩ tương đương
chỉ có hai cặp số nguyên.

Nhiệm vụ là đọc vào toàn bộ các cặp số nguyên của một kỵ sĩ, rồi tìm một kỵ
sĩ tương đương chỉ dùng đúng hai cặp số nguyên.

\section{Mô hình toán học}

Xem mỗi cặp di chuyển là một vectơ
\[
  v_i=(x_i,y_i)\in\mathbb{Z}^2,\qquad i=1,2,\ldots,n.
\]
Sau nhiều lần di chuyển, kỵ sĩ có thể đến đúng các điểm có dạng
\[
  \sum_{i=1}^{n} c_i v_i,\qquad c_i\in\mathbb{Z}.
\]
Nói cách khác, tập các điểm đến được là mô-đun con của \(\mathbb{Z}^2\) sinh
bởi các vectơ đã cho. Bài toán yêu cầu tìm hai vectơ \(u,w\) sao cho
\[
  \left\{\sum_{i=1}^{n} c_i v_i:c_i\in\mathbb{Z}\right\}
  =
  \{au+bw:a,b\in\mathbb{Z}\}.
\]

\section{Ý tưởng tổng quát}

Một cách nghĩ trực tiếp là liệt kê hai vectơ cần tìm rồi dùng tìm kiếm theo
bề rộng để kiểm tra tương đương. Cách này không khả thi: phạm vi của hai vectơ
cần tìm là vô hạn, và bàn cờ cũng vô hạn.

Các hướng tham lam, quy hoạch động hay đồ thị thông thường cũng không phù hợp.
Vì bản chất của bài toán là biến đổi tập sinh của một lưới nguyên, con đường
tự nhiên là dùng phương pháp toán học.

Nếu ba vectơ bất kỳ luôn có thể được thay bằng hai vectơ tương đương, thì với
\(n>2\) vectơ, ta chỉ cần chọn ba vectơ, thay chúng bằng hai vectơ, làm số
vectơ giảm đi \(1\), rồi lặp lại cho đến khi còn đúng hai vectơ. Vì vậy trọng
tâm là: làm thế nào để thay ba vectơ bằng hai vectơ tương đương?

\section{Ràng buộc bài toán con}

Điều kiện ``tìm hai vectơ bất kỳ tương đương'' còn quá rộng, nên khó nhìn ra
cấu trúc. Ta trước hết giải một nhiệm vụ hẹp hơn: với hai vectơ, biến chúng
thành hai vectơ tương đương có dạng đặc biệt, trong đó một vectơ là vectơ
thẳng đứng.

Xét hai vectơ không cùng phương
\[
  p=(a,b),\qquad q=(c,d).
\]
Đặt
\[
  g=\gcd(a,c),\qquad \Delta=ad-bc.
\]
Nếu một trong hai vectơ đã thẳng đứng, ta có thể giữ nguyên cặp vectơ. Ngược
lại, nhờ thuật toán Euclid mở rộng, tìm được hai số nguyên \(s,t\) sao cho
\[
  sa+tc=g.
\]
Khi đó vectơ
\[
  r=sp+tq=(g,sb+td)
\]
nằm trong lưới sinh bởi \(p,q\) và có hoành độ nhỏ nhất theo nghĩa chia hết.
Mặt khác, các điểm trên cùng một đường thẳng đứng trong lưới xuất hiện cách
nhau một khoảng
\[
  h=\frac{|\Delta|}{g}.
\]
Vì vậy cặp
\[
  r=(g,sb+td),\qquad z=(0,h)
\]
sinh đúng cùng một lưới với \(p,q\).

\section{Tính chất của hai vectơ}

Nguồn gốc minh họa bằng ví dụ \((2,3)\) và \((4,2)\). Từ hình vẽ có thể quan
sát các tính chất sau:

\begin{enumerate}
  \item Các điểm đến được nằm trên các đường thẳng đứng có hoành độ là bội
  của \(g=\gcd(a,c)\), và chỉ trên các đường đó.
  \item Trên trục \(Oy\), các điểm đến được xuất hiện tuần hoàn theo bước
  \(h=|\Delta|/g\).
  \item Các điểm trên mọi đường thẳng đứng nói trên có thể thu được bằng cách
  tịnh tiến các điểm trên trục \(Oy\) bởi vectơ \(r\).
\end{enumerate}

Chứng minh tính chất thứ hai là phần then chốt. Với mọi tổ hợp nguyên
\(\alpha p+\beta q\) có hoành độ \(0\), ta có
\[
  \alpha a+\beta c=0.
\]
Do \(g=\gcd(a,c)\), tồn tại một số nguyên \(m\) sao cho
\[
  \alpha=m\frac{c}{g},\qquad
  \beta=-m\frac{a}{g}.
\]
Tung độ tương ứng là
\[
  \alpha b+\beta d
  =
  m\frac{cb-ad}{g}
  =
  -m\frac{\Delta}{g}.
\]
Vì vậy trên trục \(Oy\), mọi điểm có tung độ là bội của \(h\), và ngược lại
các bội đó đều đạt được. Điều này giải thích vì sao vectơ thẳng đứng \(z\) có
thể thay thế toàn bộ phần dao động theo phương dọc.

\section{Chứng minh phép thay thế hai vectơ}

Ta cần chứng minh hai cặp \(\{p,q\}\) và \(\{r,z\}\) sinh cùng một lưới.

Trước hết, \(r=sp+tq\) nên \(r\) thuộc lưới sinh bởi \(p,q\). Vectơ \(z\) cũng
thuộc lưới đó vì nó là một tổ hợp nguyên tạo hoành độ \(0\) và tung độ
\(\pm h\). Do đó mọi tổ hợp của \(r,z\) đều đạt được bằng \(p,q\).

Ngược lại, xét một điểm bất kỳ
\[
  X=\alpha p+\beta q.
\]
Hoành độ của \(X\) là \(\alpha a+\beta c\), nên là bội của \(g\). Giả sử
hoành độ đó bằng \(\lambda g\). Vì \(r\) có hoành độ \(g\), điểm
\[
  X-\lambda r
\]
nằm trên trục \(Oy\). Theo tính chất trên, nó là một bội nguyên của \(z\).
Vậy \(X\) cũng là tổ hợp nguyên của \(r,z\). Hai cặp vectơ vì thế tương đương.

\section{Thay ba vectơ bằng hai vectơ}

Bây giờ quay lại nhiệm vụ chính. Với ba vectơ \(v_1,v_2,v_3\), làm như sau:

\begin{enumerate}
  \item Dùng phép thay thế hai vectơ để biến \(v_1,v_2\) thành một cặp tương
  đương gồm một vectơ thẳng đứng \(z_1\) và một vectơ còn lại \(r_1\).
  \item Dùng tiếp phép thay thế hai vectơ để biến \(z_1,v_3\) thành một cặp
  tương đương gồm một vectơ thẳng đứng \(z_2\) và một vectơ còn lại \(r_2\).
  \item Khi đó ba vectơ ban đầu tương đương với \(r_1,z_2,r_2\). Nếu một trong
  các vectơ còn lại đã sinh được bởi hai vectơ kia thì bỏ nó. Nếu chưa, áp
  dụng lại phép chuẩn hóa cho hai vectơ không thẳng đứng để thu được một cặp
  sinh tương đương.
\end{enumerate}

Diễn giải theo đại số tuyến tính nguyên, các bước trên chính là liên tục dùng
thuật toán Euclid mở rộng để thay đổi cơ sở của lưới mà không làm thay đổi tập
các điểm sinh ra. Mỗi lần xử lý ba vectơ, ta giảm số vectơ đi một.

\section{Thuật toán hoàn chỉnh}

Bắt đầu với danh sách \(n\) vectơ đã cho. Khi danh sách còn nhiều hơn hai
vectơ, chọn ba vectơ bất kỳ trong danh sách, thay chúng bằng hai vectơ tương
đương theo quy trình ở trên. Lặp lại cho đến khi danh sách còn đúng hai vectơ.
Hai vectơ cuối cùng chính là đáp án.

Mỗi lần thay thế chỉ cần các phép tính \(\gcd\) và Euclid mở rộng, nên chi phí
là \(O(\log M)\), với \(M\) là độ lớn tọa độ đang xét. Vì số lần thay thế là
\(O(n)\), tổng thời gian là
\[
  O(n\log M),
\]
và bộ nhớ phụ là \(O(1)\), ngoài danh sách các vectơ hiện tại.

\section{Kết luận}

Bản chất của tư tưởng ``thêm ràng buộc'' trong bài này là thu hẹp phạm vi tìm
kiếm nhưng vẫn bảo đảm còn nghiệm. Khi xác định thuật toán, ta ràng buộc quá
trình phải là ``biến ba thành hai''. Khi biến đổi hai vectơ, ta ràng buộc một
vectơ kết quả phải thẳng đứng. Khi chứng minh, ta tiếp tục cố định một số biến
để làm lộ ra chu kỳ theo trục \(Oy\). Nhờ các ràng buộc hợp lý này, bài toán
từ một không gian vô hạn trở thành một chuỗi phép tính \(\gcd\) đơn giản.

\end{document}
